% TFG.tex
% TFG
% Autor
% Sevilla, 16 de julio de 2016
% =============================================================================

\documentclass[a4paper,12pt,twoside]{book}

%%%%%%%%%%%%%%%%%%%%%%%%%%%%%%%%%%%%%%%%%%%%%%%%%%%%%%%%%%%%%%%%%%%%%%%%%%%%%%
%%  Paquetes adicionales                                                   %%
%%%%%%%%%%%%%%%%%%%%%%%%%%%%%%%%%%%%%%%%%%%%%%%%%%%%%%%%%%%%%%%%%%%%%%%%%%%%%%

\usepackage[utf8x]{inputenc}       % Acentos de UTF8
\usepackage[spanish]{babel}        % Castellanización.
\usepackage[T1]{fontenc}           % Codificación T1 con European Computer
                                   % Modern.  
\usepackage{graphicx}
\usepackage{fancyvrb}              % Verbatim extendido
\usepackage{makeidx}               % Índice
\usepackage{amsmath}               % AMS LaTeX
\usepackage{amsthm} 
\usepackage{latexsym}
\usepackage[colorinlistoftodos
           , backgroundcolor = yellow
           , textwidth = 4cm
           , shadow
           , spanish]{todonotes}
% Fuentes
\usepackage{mathpazo}              % Fuentes semejante a palatino
\usepackage[scaled=.90]{helvet}
\usepackage{cmtt}
\renewcommand{\ttdefault}{cmtt}
\usepackage{a4wide}
% \usepackage{xmpincl}               % Incluye metadato de licencia CC.

% Tikz
\usepackage{tkz-berge}
\usetikzlibrary{shapes,trees}

% Desactivar <,> cuando se hacen dibujos con tikz.
\tikzset{
every picture/.append style={
  execute at begin picture={\deactivatequoting},
  execute at end picture={\activatequoting}
  }
}

% Márgenes
\usepackage[margin=1in]{geometry}

% Control de espacios en la tabla de contenidos:
\usepackage[titles]{tocloft}
\setlength{\cftbeforechapskip}{2ex}
\setlength{\cftbeforesecskip}{0.5ex}
\setlength{\cftsecnumwidth}{12mm}
\setlength{\cftsubsecindent}{18mm}

% Control de listas
% Elimina espacio entre item y párrafo y coloca item en el margen izquierdo
% \usepackage{enumitem}
% \setlist[enumerate,itemize]{noitemsep, nolistsep, leftmargin=*}

\usepackage{minitoc}

% Doble espacio entre líneas
\usepackage{setspace}
\onehalfspacing


% \linespread{1.05}                % Distancia entre líneas
\setlength{\parindent}{5mm}        % Indentación de comienzo de párrafo
\deactivatetilden                  % Elima uso de ~ para la eñe
\raggedbottom                      % No ajusta los espacios verticales.

\usepackage[%
  pdftex,
  pdfauthor={Eduardo Paluzo},%
  pdftitle={LPO en Haskell},%
  pdfstartview=FitH,%
  bookmarks=false,%
  colorlinks=true,%
  urlcolor=blue,%
  unicode=true]{hyperref}      

\usepackage{tikz}

%%%%%%%%%%%%%%%%%%%%%%%%%%%%%%%%%%%%%%%%%%%%%%%%%%%%%%%%%%%%%%%%%%%%%%%%%%%%%%
%%  Cabeceras                                                              %%
%%%%%%%%%%%%%%%%%%%%%%%%%%%%%%%%%%%%%%%%%%%%%%%%%%%%%%%%%%%%%%%%%%%%%%%%%%%%%%

\usepackage{fancyhdr}

\addtolength{\headheight}{\baselineskip}

\pagestyle{fancy}

\cfoot{}

\fancyhead{}
\fancyhead[RE]{\bfseries \nouppercase{\leftmark}}
\fancyhead[LO]{\bfseries \nouppercase{\rightmark}}
\fancyhead[LE,RO]{\bfseries \thepage}

%%%%%%%%%%%%%%%%%%%%%%%%%%%%%%%%%%%%%%%%%%%%%%%%%%%%%%%%%%%%%%%%%%%%%%%%%%%%%%
%%  Definiciones                                                           %%
%%%%%%%%%%%%%%%%%%%%%%%%%%%%%%%%%%%%%%%%%%%%%%%%%%%%%%%%%%%%%%%%%%%%%%%%%%%%%%

\input definiciones
\def\mtctitle{Contenido}

%%%%%%%%%%%%%%%%%%%%%%%%%%%%%%%%%%%%%%%%%%%%%%%%%%%%%%%%%%%%%%%%%%%%%%%%%%%%%%
%%  Título                                                                 %%
%%%%%%%%%%%%%%%%%%%%%%%%%%%%%%%%%%%%%%%%%%%%%%%%%%%%%%%%%%%%%%%%%%%%%%%%%%%%%%

\title{\Huge Título del TFG}
\author{Autor}
\date{\vfill \hrule \vspace*{2mm}
  \begin{tabular}{l}
      \href{http://www.cs.us.es/glc}
           {Grupo de Lógica Computacional} \\
      \href{http://www.cs.us.es}
           {Dpto. de Ciencias de la Computación e Inteligencia Artificial} \\
      \href{http://www.us.es}
           {Universidad de Sevilla}  \\
      Sevilla, 16 de junio de 2016 (Versión de \today)
  \end{tabular}\hfill\mbox{}}




%%%%%%%%%%%%%%%%%%%%%%%%%%%%%%%%%%%%%%%%%%%%%%%%%%%%%%%%%%%%%%%%%%%%%%%%%%%%%%%
%%  Construcción del índice                                                 %%
%%%%%%%%%%%%%%%%%%%%%%%%%%%%%%%%%%%%%%%%%%%%%%%%%%%%%%%%%%%%%%%%%%%%%%%%%%%%%%%

\makeindex

%%%%%%%%%%%%%%%%%%%%%%%%%%%%%%%%%%%%%%%%%%%%%%%%%%%%%%%%%%%%%%%%%%%%%%%%%%%%%%
%%  Documento                                                              %%
%%%%%%%%%%%%%%%%%%%%%%%%%%%%%%%%%%%%%%%%%%%%%%%%%%%%%%%%%%%%%%%%%%%%%%%%%%%%%%

% \includeonly{Introduccion}

% \includexmp{licencia}

\begin{document}

\dominitoc

\begin{titlepage}
 \vspace*{2cm}
  \begin{center}
    {\huge \textbf{Álgebra constructiva en Haskell}}
  \end{center}
  \vspace{4cm}
  \begin{center}
    \leavevmode\includegraphics[totalheight=6cm]{sello.png}\\[3cm]
    {\normalsize Facultad de Matemáticas} \\
    {\normalsize Departamento de Ciencias de la Computación e Inteligencia Artificial}\\
    {\normalsize Trabajo Fin de Grado} \\
  \end{center}
  \begin{center}
    {\large \textbf{Autor}}
  \end{center}
  \newpage
 
 \begin{flushright}
   \vspace*{5cm}
   \begin{minipage}{8.45cm}
      Agradecimientos
    \end{minipage}

      \vspace*{7.5mm}

  \end{flushright}
  \vspace*{\fill}

  \newpage


  
  % \begin{flushright}
  \begin{center}
   \vspace*{5cm}
    \begin{minipage}{14cm}
      El presente Trabajo Fin de Grado se ha realizado en el Departamento de
      Ciencias de la Computación e Inteligencia Artificial de la Universidad de
      Sevilla.

      \vspace*{7.5mm}

      Supervisado por
      % \vspace*{5mm}
    \end{minipage}\par
    Tutor
    % \end{flushright}
    \end{center}
  \vspace*{\fill}

  \newpage

  \vspace*{3cm}
  {\huge \textit{Abstract}}

  \vspace{2cm}
  Resumen en inglés
  
\end{titlepage}
\newpage

\noindent
Esta obra está bajo una licencia Reconocimiento--NoComercial--CompartirIgual
2.5 Spain de Creative Commons. 

\vspace*{4ex}

\begin{center}
\fbox{
\begin{minipage}{175mm}
\vspace*{1ex}
\noindent                     
{\bf Se permite:}
\begin{itemize}                         
\item copiar, distribuir y comunicar públicamente la obra
\item hacer obras derivadas
\end{itemize}
         
\noindent       
{\bf Bajo las condiciones siguientes:}

\begin{minipage}[t]{20mm}
  \vspace*{0mm}\hspace*{1em}
  \includegraphics[scale=0.35]{licencia/deed.png}
\end{minipage}%
\begin{minipage}[t]{15cm}
  \vspace*{0mm}
  {\bf Reconocimiento}. Debe reconocer los créditos de la  obra de la manera
  especificada por el autor.  
\end{minipage}
\\[1ex]

\begin{minipage}[t]{20mm}
  \vspace*{0mm}\hspace*{1em}
  \includegraphics[scale=0.35]{licencia/deed-eu.png}
\end{minipage}%
\begin{minipage}[t]{15cm}
  \vspace*{2mm}
  {\bf No comercial}. No puede utilizar esta obra para fines comerciales.
\end{minipage}
\\[1ex]

\begin{minipage}[t]{20mm}
  \vspace*{0mm}\hspace*{1em}
  \includegraphics[scale=0.35]{licencia/deed_002.png}
\end{minipage}%
\begin{minipage}[t]{15cm}
  \vspace*{-2mm}
  {\bf Compartir bajo la misma licencia}. Si altera o transforma esta obra, o
  genera una obra derivada, sólo puede distribuir la obra generada bajo una
  licencia idéntica a ésta. 
\end{minipage}

\begin{itemize}    
\item Al reutilizar o distribuir la obra, tiene que dejar bien claro los
  términos de la licencia de esta obra. 
\item Alguna de estas condiciones puede no aplicarse si se obtiene el permiso
 del titular de los derechos de autor.
\end{itemize}
\vspace*{1ex}
\end{minipage}}
\end{center}

\vspace*{4ex}

Esto es un resumen del texto legal (la licencia completa). Para ver una copia
de esta licencia, visite 
\href{http://creativecommons.org/licenses/by-nc-sa/2.5/es/}
     {\url{http://creativecommons.org/licenses/by-nc-sa/2.5/es/}}
o envie una carta a Creative Commons, 559 Nathan Abbott Way, Stanford,
California 94305, USA.

%%% Local Variables: 
%%% mode: latex
%%% TeX-master: t
%%% End: 

\newpage 

\tableofcontents
\newpage

\chapter{Programación funcional con Haskell}\label{sec:progfunHas}

En este capítulo se hace una breve introducción a la programación funcional en
Haskell suficiente para entender su aplicación en los siguientes
capítulos. Para una introducción más amplia se pueden consultar los apuntes de
la asignatura de Informática de 1º del Grado en Matemáticas
(\cite{Alonso-15b}). 

El contenido de este capítulo se encuentra en el módulo \texttt{PFH} 
\entrada{PFH}.

%%% Local Variables:
%%% mode: latex
%%% TeX-master: "TFG"
%%% End:


%%%%%%%%%%%%%%%%%%%%%%%%%%%%%%%%%%%%%%%%%%%%%%%%%%%%%%%%%%%%%%%%%%%%%%%%%%%%%%% 
%%  Bibliografía                                                            %%
%%%%%%%%%%%%%%%%%%%%%%%%%%%%%%%%%%%%%%%%%%%%%%%%%%%%%%%%%%%%%%%%%%%%%%%%%%%%%%%

\addcontentsline{toc}{chapter}{Bibliografía}
\bibliographystyle{abbrv}
\bibliography{TFG}

%%%%%%%%%%%%%%%%%%%%%%%%%%%%%%%%%%%%%%%%%%%%%%%%%%%%%%%%%%%%%%%%%%%%%%%%%%%%%%%
%%  Índice                                                                  %%
%%%%%%%%%%%%%%%%%%%%%%%%%%%%%%%%%%%%%%%%%%%%%%%%%%%%%%%%%%%%%%%%%%%%%%%%%%%%%%%

\addcontentsline{toc}{chapter}{Indice de definiciones}

\printindex

%%%%%%%%%%%%%%%%%%%%%%%%%%%%%%%%%%%%%%%%%%%%%%%%%%%%%%%%%%%%%%%%%%%%%%%%%%%%%%%
%% § Pendientes                                                              %%
%%%%%%%%%%%%%%%%%%%%%%%%%%%%%%%%%%%%%%%%%%%%%%%%%%%%%%%%%%%%%%%%%%%%%%%%%%%%%%%

%\todototoc
%\listoftodos

\end{document}



%%% Local Variables:
%%% mode: latex
%%% TeX-master: t
%%% End:
